\documentclass{article}
\usepackage[utf8]{inputenc}
\usepackage{geometry}
\geometry{a4paper, margin=1in}
\usepackage{hyperref}
\usepackage{listings}
\usepackage{xcolor}

\definecolor{codegreen}{rgb}{0,0.6,0}
\definecolor{codegray}{rgb}{0.5,0.5,0.5}
\definecolor{codepurple}{rgb}{0.58,0,0.82}
\definecolor{backcolour}{rgb}{0.95,0.95,0.92}

\lstdefinestyle{mystyle}{
    backgroundcolor=\color{backcolour},   
    commentstyle=\color{codegreen},
    keywordstyle=\color{magenta},
    numberstyle=\tiny\color{codegray},
    stringstyle=\color{codepurple},
    basicstyle=\ttfamily\footnotesize,
    breakatwhitespace=false,         
    breaklines=true,                 
    captionpos=b,                    
    keepspaces=true,                 
    numbers=left,                    
    numbersep=5pt,                  
    showspaces=false,                
    showstringspaces=false,
    showtabs=false,                  
    tabsize=2
}
\lstset{style=mystyle}

\title{Secure Password Manager: Vault Export Walkthrough}
\author{Documentation}
\date{\today}

\begin{document}

\maketitle

\section{Overview}
The \texttt{export.py} module implements a secure protocol to transfer a user's entire password vault to another user. Since vaults are securely encrypted and tied to a user's master password, we cannot simply copy the file. Instead, the module uses a combination of asymmetric cryptography (Diffie-Hellman Key Exchange, ElGamal Digital Signatures) and symmetric cryptography (AES-256-GCM) to ensure that the transfer is secure against eavesdropping and tampering.

\section{The Export Protocol Walkthrough}

The vault export process is divided into three distinct phases:

\subsection{Phase 1: Diffie-Hellman Key Exchange}
To securely encrypt the vault for transit, both devices need a shared secret key that an attacker cannot intercept.
\begin{enumerate}
    \item Both the Sender (Device 1) and Receiver (Device 2) generate ephemeral (temporary) Diffie-Hellman public and private keys using shared global parameters ($q$ and $\alpha$).
    \item Device 1 signs its ephemeral public key using its ElGamal private key and sends it to Device 2.
    \item Device 2 verifies Device 1's signature. This prevents Man-in-the-Middle (MitM) attacks.
    \item Device 2 does the same: signs its public key and sends it to Device 1, which verifies it.
    \item Both devices compute the shared DH secret and derive a 256-bit AES session key using SHA-256.
\end{enumerate}

\subsection{Phase 2: Vault Transfer (Sender Side)}
With a secure session key established, the Sender prepares the data.
\begin{enumerate}
    \item The Sender decrypts their local vault using their master password.
    \item The Sender re-encrypts the vault data using the AES-256-GCM session key derived in Phase 1.
    \item The Sender signs the encrypted vault ciphertext using their ElGamal private key to guarantee authenticity.
    \item The encrypted data, cryptographic nonce, authentication tag, and digital signature are packaged and ``transmitted'' to the Receiver.
\end{enumerate}

\subsection{Phase 3: Vault Import (Receiver Side)}
The Receiver processes the incoming secure package.
\begin{enumerate}
    \item The Receiver verifies the digital signature of the package using the Sender's public ElGamal key. If invalid, the process aborts.
    \item The Receiver decrypts the vault using the AES-256-GCM session key.
    \item The decrypted credentials are saved into the Receiver's local vault, encrypted with the Receiver's master password.
\end{enumerate}

\section{Function Explanations}

Here is a breakdown of the helper functions defined in \texttt{export.py}:

\begin{itemize}
    \item \textbf{\texttt{load\_dh\_params()}}: Reads the Diffie-Hellman parameters (a large prime $q$ and a primitive root $\alpha$) from \texttt{config/params.json}.
    \item \textbf{\texttt{generate\_dh\_keypair(q, alpha)}}: Generates a random private key and computes the corresponding public key $\alpha^{\text{private}} \pmod q$.
    \item \textbf{\texttt{compute\_shared\_secret(other\_pub, my\_priv, q)}}: Computes the Diffie-Hellman shared secret: $\text{other\_pub}^{\text{my\_priv}} \pmod q$.
    \item \textbf{\texttt{derive\_session\_key(shared\_secret)}}: Hashes the large integer shared secret using SHA-256 to produce a secure 256-bit (32-byte) key for AES encryption.
    \item \textbf{\texttt{\_encrypt\_with\_key(plaintext\_bytes, key)}}: Encrypts data using AES-GCM, which provides both confidentiality and authenticated encryption (preventing tampering).
    \item \textbf{\texttt{\_decrypt\_with\_key(ciphertext, nonce, tag, key)}}: Decrypts AES-GCM data, verifying the authentication tag to ensure the ciphertext was not modified.
    \item \textbf{\texttt{export\_vault(...)}}: The main orchestration function that runs Phase 1, Phase 2, and Phase 3 in sequence.
\end{itemize}

\section{Line-by-Line Code Breakdown}

Below is the complete \texttt{export.py} code with line-by-line explanations.

\begin{lstlisting}[language=Python]
#mod4
import json            # For parsing JSON data and config files
import os              # For filesystem paths
import secrets         # For cryptographically secure random number generation
import hashlib         # For SHA-256 hashing
import base64          # For encoding binary ciphertext to strings
from Crypto.Cipher import AES  # PyCryptodome for AES-GCM encryption
from modules import elgamal    # Local module: Key management
from modules import signatures # Local module: Digital signatures
from modules import vault      # Local module: Vault CRUD operations

# Define paths to locate the config/params.json file dynamically
MODULES_DIR = os.path.dirname(os.path.abspath(__file__))
PROJECT_ROOT = os.path.dirname(MODULES_DIR)
CONFIG_PATH = os.path.join(PROJECT_ROOT, "config", "params.json")

def load_dh_params():
    # Open the params.json file in read mode
    with open(CONFIG_PATH, 'r') as f:
        data = json.load(f) # Parse JSON to a Python dictionary
        # Extract DH parameters 'q' and 'alpha', converting from hex strings to integers
        q = int(data['dh']['q'], 16)
        alpha = int(data['dh']['alpha'], 16)
        return q, alpha # Return the parameters tuple

def generate_dh_keypair(q, alpha):
    # Generate a random private key between 2 and q-2 using secure RNG
    private_key = 2 + secrets.randbelow(q - 3)
    # Calculate public key = (alpha^private_key) mod q
    public_key = pow(alpha, private_key, q)
    return public_key, private_key # Return the keypair tuple

def compute_shared_secret(other_public_key, my_private_key, q):
    # Calculate shared secret = (other_public_key^my_private_key) mod q
    return pow(other_public_key, my_private_key, q)

def derive_session_key(shared_secret):
    # Convert the large integer shared secret into big-endian bytes
    secret_bytes = shared_secret.to_bytes(
        (shared_secret.bit_length() + 7) // 8, byteorder='big'
    )
    # Hash the bytes with SHA-256 to create a uniformly random 256-bit key
    return hashlib.sha256(secret_bytes).digest()

def _encrypt_with_key(plaintext_bytes, key):
    # Initialize AES cipher in Galois/Counter Mode (GCM) using the session key
    cipher = AES.new(key, AES.MODE_GCM)
    # Encrypt the plaintext and generate an authentication tag
    ciphertext, tag = cipher.encrypt_and_digest(plaintext_bytes)
    # Return the encrypted data, the unique nonce used, and the auth tag
    return ciphertext, cipher.nonce, tag

def _decrypt_with_key(ciphertext, nonce, tag, key):
    # Initialize AES cipher in GCM mode with the provided key and nonce
    cipher = AES.new(key, AES.MODE_GCM, nonce=nonce)
    # Decrypt the ciphertext and verify the authentication tag matches
    return cipher.decrypt_and_verify(ciphertext, tag)

def export_vault(sender_username, sender_master_pw,
                 receiver_username, receiver_master_pw):
    # Load ElGamal private keys for both users into memory
    elgamal.load_private_key(sender_username)
    elgamal.load_private_key(receiver_username)

    print("\n[Phase 1] Diffie-Hellman Key Exchange")
    # Load global DH parameters q and alpha
    q, alpha = load_dh_params()
    # Sender generates ephemeral DH keys
    sender_dh_pub, sender_dh_priv = generate_dh_keypair(q, alpha)
    # Receiver generates ephemeral DH keys
    receiver_dh_pub, receiver_dh_priv = generate_dh_keypair(q, alpha)
    print("  Both devices generated ephemeral DH key pairs.")

    # Sender creates an ElGamal signature of their DH public key
    sig1 = signatures.sign_vault(str(sender_dh_pub), sender_username)
    print(f"  Device 1 ({sender_username}) signed and sent DH public key.")

    # Receiver verifies the Sender's signature to prevent MitM
    signatures.verify_vault(str(sender_dh_pub), sig1, sender_username)
    print("  Device 2 verified Device 1's DH public key signature. [OK]")

    # Receiver creates an ElGamal signature of their DH public key
    sig2 = signatures.sign_vault(str(receiver_dh_pub), receiver_username)
    print(f"  Device 2 ({receiver_username}) signed and sent DH public key.")

    # Sender verifies the Receiver's signature
    signatures.verify_vault(str(receiver_dh_pub), sig2, receiver_username)
    print("  Device 1 verified Device 2's DH public key signature. [OK]")

    # Sender computes the shared secret using Receiver's pubkey and their own privkey
    sender_shared = compute_shared_secret(receiver_dh_pub, sender_dh_priv, q)
    # Receiver computes the shared secret using Sender's pubkey and their own privkey
    receiver_shared = compute_shared_secret(sender_dh_pub, receiver_dh_priv, q)
    # Sanity check: verify the math worked and both secrets are identical
    assert sender_shared == receiver_shared, "Shared secrets do not match!"
    print("  Shared secret computed on both devices. [OK]")

    # Derive the 256-bit AES session key from the shared secret
    session_key = derive_session_key(sender_shared)
    print("  Session AES-256 key derived via SHA-256. [OK]")
    
    print("\n[Phase 2] Vault Transfer")
    # Sender decrypts their own vault using their master password
    creds = vault.load_vault(sender_username, sender_master_pw)
    print(f"  Device 1 decrypted vault ({len(creds)} credential(s)).")

    # Serialize the credentials list to a JSON string, then encode to bytes
    plain_json = json.dumps(creds).encode('utf-8')
    # Encrypt the vault bytes using AES-GCM and the DH session key
    session_ct, session_nonce, session_tag = _encrypt_with_key(
        plain_json, session_key
    )
    # Base64 encode the binary outputs so they can be transmitted as text
    session_ct_b64 = base64.b64encode(session_ct).decode()
    session_nonce_b64 = base64.b64encode(session_nonce).decode()
    session_tag_b64 = base64.b64encode(session_tag).decode()
    print("  Vault re-encrypted with DH session key.")

    # Sender creates an ElGamal signature over the encrypted vault data
    transfer_sig = signatures.sign_vault(session_ct_b64, sender_username)
    print("  Session-encrypted data signed by Device 1. [OK]")

    # Package the encrypted data, AES parameters, and ElGamal signature
    transfer_package = {
        "encrypted_data": session_ct_b64,
        "nonce": session_nonce_b64,
        "tag": session_tag_b64,
        "signature": {"r": transfer_sig['r'], "s": transfer_sig['s']},
        "sender": sender_username,
    }

    print("\n[Phase 3] Vault Import")
    # Receiver verifies the Sender's ElGamal signature on the transfer package
    signatures.verify_vault(
        transfer_package["encrypted_data"],
        transfer_package["signature"],
        sender_username
    )
    print("  Device 2 verified transfer signature. [OK]")

    # Receiver base64-decodes the incoming data back into bytes
    received_ct = base64.b64decode(transfer_package["encrypted_data"])
    received_nonce = base64.b64decode(transfer_package["nonce"])
    received_tag = base64.b64decode(transfer_package["tag"])

    # Receiver decrypts the vault using the AES-GCM session key
    decrypted_json = _decrypt_with_key(
        received_ct, received_nonce, received_tag, session_key
    )

    # Parse the decrypted JSON bytes back into a Python list of dictionaries
    imported_creds = json.loads(decrypted_json.decode('utf-8'))
    print(f"  Decrypted {len(imported_creds)} credential(s) with session key.")
    
    # Receiver saves the imported credentials into their own vault, encrypted with their password
    vault.save_vault(receiver_username, receiver_master_pw, imported_creds)
    print(f"  Vault saved for '{receiver_username}' with new master password.")

    # Print success message
    print(f"\n[OK] Export complete! {len(imported_creds)} credential(s) "
          f"transferred from '{sender_username}' to '{receiver_username}'.")
    # Return the imported credentials
    return imported_creds
\end{lstlisting}

\end{document}
